% Member 4 - final report, III.F The Strength of the Style as the Controlled Variable
\subsection{The Strength of the Style as the Controlled Variable}
\label{sec:strength_energy}

The central quantity of this study is not the spatial distribution of the
style but its \emph{amount}. This section defines that amount precisely and
explains why it has to be controlled before any claim about the
distribution can be made.

\paragraph{Two different notions of strength.}
A stylisation has at least two natural measures of how much it changed the
image, and they are not the same thing. The first is the \emph{nominal}
strength $\alpha_0$, the parameter that appears in the blending rule. The
second is the \emph{effective} strength $\bar{\alpha}$ of
Equation~\eqref{eq:meanstrength}, the average of the blending field over
the feature grid. For a uniform configuration the two coincide, because
$\alpha$ is constant. For a modulated configuration they do not: the field
of Equation~\eqref{eq:alphamap} averages
$\alpha_0(1-\beta\bar{M})$, which at the main setting
($\alpha_0 = 0.5$, $\beta = 0.7$, measured $\bar{M} = 0.361$) gives
$\bar{\alpha} = 0.374$ rather than $0.5$.

\paragraph{Why the difference is not negligible.}
The difference between $0.5$ and $0.374$ is large on the scale of the
effect we are trying to detect: it is exactly the difference between the
uniform AdaIN group and the strength-matched control, and those two differ
by about two accuracy points. A comparison between a modulated field and a
uniform style at the same $\alpha_0$ therefore measures two things at once,
and the two do not have the same sign: the modulation is expected to help
(the claim under test), while the reduced effective strength is expected to
hurt (a well-known effect of weaker augmentation). Any result of such a
comparison is uninterpretable.

\paragraph{How the effective strength is measured.}
We do not assume the relation $\bar{\alpha} = \alpha_0(1-\beta\bar{M})$;
we measure it. For every generated image, the actual blending field is
averaged over the feature grid and recorded, and the effective strength of
the whole set is the mean of those per-image values. Over the 1{,}349
training images of the main configuration the measured value is
$0.3736 \pm 0.0273$, in agreement with the nominal prediction of $0.374$.
The measurement is what the strength-matched control is built from: its
constant $\alpha$ is set to the measured mean of the field it has to match,
not to the nominal $\alpha_0$.

\paragraph{The two controls.}
We use two matched controls, one for each direction of the comparison, and
they answer different questions.

The \emph{strength-matched uniform AdaIN} replaces the modulated field by
the constant that has the same average strength. Comparing against it
answers the question we actually care about: \emph{given the same total
amount of style, does it matter how that amount is distributed over the
image?} The answer we obtain is no.

The \emph{equal-strength spatial map} goes the other way: it keeps the
modulation and recalibrates $\alpha_0$ so that the mean strength equals
that of a uniform AdaIN at $\alpha = 0.5$, namely
$\alpha_0 = 0.500/(1-0.7\times0.361) = 0.6692$. Comparing against it
answers the complementary question: \emph{at a fixed nominal strength, how
much does the modulation change the result, once the change in total
strength is compensated?} Here the two groups differ in the modulation and
in nothing else.

Reporting both is what makes the conclusion robust to the choice of which
quantity is held fixed. If the spatial distribution mattered, then
\emph{both} comparisons would show an advantage for the modulated field,
because both remove the strength confound. They do not. The advantage
appears only when the strength is \emph{not} controlled, which is precisely
the confound, not an effect.
